\section{DB、DC間の移行等}

\subsection{過去分の一部DC移行}

\begin{InyouJobunbox}
  \begin{itemize}
    \item 法第82条の2
    \item 令第54条の2、第54条の4
    \item 則第96条の2、第96条の5〜第96条の6
  \end{itemize}
\end{InyouJobunbox}

\begin{fillblank}[年金2 2018 問題1(3)B,C]
  \noindent \textbf{◯DB令}
  \begin{lawstructure}
    \begin{lawarticle}{第五十四条の二}{確定拠出年金を実施する場合の積立金の移換}
      法第八十二条の二第一項の規定による積立金の移換は、次に定めるところにより行うものとする。

      \lawitem{一}
      加入者の給付の額を減額することにより当該加入者の個人別管理資産（確定拠出年金法第二条第十二項に規定する個人別管理資産をいう。以下同じ。）に充てるものであること。

       \lawitem{二}
      移換加入者（法第八十二条の二第二項に規定する移換加入者をいう。以下同じ。）となるべき者の範囲が同条第一項の規約において定められていること。

       \lawitem{三}
      前号の移換加入者となるべき者の範囲は、特定の者について不当に差別的なものでなく、かつ、加入者が（　　A　　）できるものでないこと。
      
      （略）    
    \end{lawarticle}

    \begin{lawarticle}{第五十四条の四}{資産の移換をする場合の掛金の一括拠出}
      事業主等が法第八十二条の二第一項の規定に基づき積立金を移換する場合において、（　　B　　）の前日における積立金のうち当該移換に係る分として厚生労働省令で定める方法により算定した額が移換加入者に係る確定拠出年金対象移換相当額の合計額を下回るときは、法第五十五条第一項の規定にかかわらず、当該移換に係る事業主は、当該下回る額を、掛金として一括して拠出しなければならない。
    \end{lawarticle}
  \end{lawstructure}
\end{fillblank}

\begin{answerbox}
  \begin{enumerate}[A: ]
    \item 任意に選択
    \item 規約変更日
  \end{enumerate}
\end{answerbox}

\begin{truefalse}[年金1 2019 問題1(2)\textcircled{2}]
  下線部の正誤を判定し、誤っている場合は正しい内容に改めなさい。
\\

  1 つの実施事業所からなる確定給付企業年金の積立金の一部を、当該実施事業所の事業主が
実施する企業型年金における加入者の個人別管理資産に移換することができるように規約に
定める場合に必要な同意は、加入者のうち移換加入者となるべき者の2 分の1 以上の同意お
よび加入者のうち移換加入者以外の加入者の\underline{3 分の2 以上の同意}である。
\end{truefalse}

\begin{answerbox}
  \textbf{誤}

  二分の一以上の同意（ならびに企業型年金を実施する事業主の全部の同意）
\end{answerbox}

\begin{truefalse}[年金1 2018 問題1(2)\textcircled{4}]
  下線部の正誤を判定し、誤っている場合は正しい内容に改めなさい。
\\

確定給付企業年金の積立金の一部を企業型年金の資産管理機関に移換する場合、企業型年金
の資産管理機関が当該資産の移換の受入れを行う日は、当該確定給付企業年金の規約変更日の
属する月の\underline{翌月末日}とする必要がある。
\end{truefalse}

\begin{answerbox}
  \textbf{誤}

  翌々日の末日以前の企業型年金規約で定める日
\end{answerbox}

\begin{explanationbox}
  DC令第22条第2項
\end{explanationbox}

\begin{truefalse}[年金1 H28 問題1(4)\textcircled{3}]
  下線部の正誤を判定し、誤っている場合は正しい内容に改めなさい。
\\

確定給付企業年金からの一部移行により、確定給付企業年金の積立金の一部を企業型年金の
資産管理機関に移換する場合、資産の移換日は、\underline{規約に定められている}必要がある。
\end{truefalse}

\begin{answerbox}
  \textbf{正}

\end{answerbox}


\begin{explanationbox}
  DC令第22条第2項
\end{explanationbox}

\begin{shortanswer}[年金1 H28 問題2(3)]
  平成 28 年 6 月に公布された「確定拠出年金法等の一部を改正する法律」の一部が平成 28 年 7
月 1 日に施行され、確定給付企業年金から確定拠出年金に積立金を移換する場合の同意要件が緩
和された。確定給付企業年金法第 82 条の 2 および確定給付企業年金法施行規則第 96 条の 5 に定
められた内容を用いて、緩和された同意要件について簡記しなさい。
\end{shortanswer}

\begin{answerbox}
  企業年金に移行しない者のみからなる実施事業所について、移行元の確定給付企業年金の掛金が増加しない場合、当該事業所の加入者の2分の 1の同意は不要となった。
\end{answerbox}

\begin{explanationbox}
  昔の時事問題
\end{explanationbox}

\begin{shortanswer}[年金2 2019 問題2(2)]
  確定給付企業年金関連法令および関連通知の内容を踏まえて、確定給付企業年金の積立金の
一部を確定拠出年金（企業型）に移換するとき、確定拠出年金（企業型）に移換しない者から
の同意が不要となる場合を 3 点簡記しなさい。
\end{shortanswer}

\begin{itemize}
  \item 移換する実施事業所が減少する場合であって、他の実施事業所の事業主の掛金が増加しない場合または当該減少事業所が掛金を一括拠出する場合
  \item 積立金の一部を移換することに伴い減少する数理債務の額から当該移換に伴い減少する特別掛金収入現価を控除した額が、当該移換に伴い減少する積立金の額を下回らない場合
  \item 移換を行う実施事業所の事業主が、移換に伴い減少する積立金の額から当該移換に伴い減少する数理債務の額を控除した額に相当する額を、過去勤務債務の額に係る特別掛金として拠出することを規約で定めている場合
\end{itemize}


\begin{explanationbox}
  類題：数理人試験 H30 問題4 設問2

類題：年金１ 2025 問題2(1)イ
\end{explanationbox}

\begin{comment}

\begin{shortanswer}[年金1 2022 問題3(2)ア 改題]
  ＤＢを制度終了し、加入者全員を対象として企業型ＤＣへ移行する場合について、次の問いに答えなさい。
  \\

上記を実施するにあたり、ＤＢ、企業型ＤＣそれぞれに必要な同意の内容について、判定根
拠も含めて簡潔に入力しなさい。なお、労働組合の同意でよいときは労働組合の同意を取得
するものとし、個別同意が必要なときは同意が必要な人数についても入力すること。
\end{shortanswer}

\begin{answerbox}
  
\end{answerbox}
  
\end{comment}

\newpage

\subsection{全部DC移行}

出題例なし（第\ajRoman{2}部にはある）

\newpage

\subsection{将来分のDC移行}

出題例なし

\newpage

\subsection{DBからDCへの脱退一時金相当額の移換}

出題例なし

\newpage

\subsection{DBから個人型DCへの残余財産の移換}


\begin{shortanswer}
  2022年５月において新たに認められた制度間の年金資産の移換（ポータビリティ）の方法を２
つ記載せよ。
\end{shortanswer}

\begin{answerbox}
  \begin{itemize}
    \item 終了した確定給付企業年金制度から個人型企業年金（iDeco）への年金資産の移換
    \item 加入者の退職等に伴う企業型確定拠出年金から通算企業年金への年金資産の移換
  \end{itemize}
\end{answerbox}

\newpage

\subsection{DBから中退共への積立金等の移換}

\begin{fillblank}[年金1 2020 問題1(2)A〜D]
  \noindent \textbf{◯DB法}
  \begin{lawstructure}
    \begin{lawarticle}{第八十二条の五}{確定給付企業年金から独立行政法人勤労者退職金共済機構への積立金等の移換}
      実施事業所の事業主が会社法（平成十七年法律第八十六号）その他の法律の規定による（　　A　　）、
      （　　B　　）その他の行為として厚生労働省令で定める行為（以下この項において「（　　A　　）等」という。）
      をした場合であって、当該（　　A　　）等に係る事業主が、当該（　　A　　）等により確定給付企業年金の加入者
      の資格を喪失した者を中小企業退職金共済法第二条第七項に規定する被共済者として同条第三項に規定する
      退職金共済契約を締結するときは、当該事業主は、政令で定める基準に従い規約で定めるところにより、
      当該加入者であった者の同意を得て、当該確定給付企業年金の（　　C　　）に機構への当該同意を得た
      加入者であった者に係る積立金（第八十三条の規定により当該確定給付企業年金が終了した場合は、
      第八十九条第六項に規定する（　　D　　））の移換を申し出ることができる。

      \lawparagraph{2}
      事業主等は、前項の規定による申出に基づき、中小企業退職金共済法第三十一条の三の規定により積立金を移換したときは、当該積立金を移換した者に係る給付の支給に関する義務を免れる。

      \lawparagraph{3}
      第一項の規定による申出に基づき、中小企業退職金共済法第三十一条の三の規定により（　　D　　）を移換したときは、
      第八十九条第六項の規定の適用については、当該（　　D　　）は、同項に規定する終了制度加入者等に分配されたものとみなす。
    \end{lawarticle}
  \end{lawstructure}
\end{fillblank}


\begin{answerbox}
  \begin{enumerate}[A: ]
    \item 合併
    \item 会社分割
    \item 資産管理運用機関等
    \item 残余財産
  \end{enumerate}
\end{answerbox}

\begin{fillblank}[年金2 2018 問題1(3)D]
  \noindent \textbf{◯DB令}
  \begin{lawstructure}
    \begin{lawarticle}{第五十四条の八}{独立行政法人勤労者退職金共済機構への積立金等の移換の基準}
      法第八十二条の五第一項の政令で定める基準は、次のとおりとする。

      \lawitem{一}
      法第八十二条の五第一項の規定による移換の申出は、同項に規定する合併等を行った日から起算して（　　A　　）を経過する日（天災その他その日までの間に申し出なかったことについてやむを得ない理由があるときは、その理由がやんだ日の属する月の翌月の末日）までの間に限って行うことができるものであること。

      （略）
     \end{lawarticle}
  \end{lawstructure}
\end{fillblank}


\begin{answerbox}
  \begin{enumerate}[A: ]
    \item 一年
  \end{enumerate}
\end{answerbox}

\newpage

\subsection{DCからDBへの資産の移換}

\begin{fillblank}[年金1 2020 問題1(2)C,E]
  \noindent \textbf{◯DB法}
  \begin{lawstructure}
    \begin{lawarticle}{第八十二条の六}{確定拠出年金又は独立行政法人勤労者退職金共済機構から確定給付企業年金への資産の移換}
      事業主等は、その（　　A　　）が確定拠出年金法第五十四条の四第二項若しくは第七十四条の四第二項の規定によりこれらの項に規定する（　　B　　）の移換を受けた場合又は中小企業退職金共済法第十七条第一項若しくは第三十一条の四第一項の規定により機構から同法第十七条第一項に規定する厚生労働省令で定める金額の引渡し若しくは同法第三十一条の四第一項に規定する解約手当金に相当する額の移換を受けた場合は、これらの金額を原資として、政令で定める基準に従い規約で定めるところにより、当該加入者に対し、老齢給付金等の支給を行うものとする。

      \lawparagraph{2}
      事業主等は、前項の規定により老齢給付金等の支給を行うこととなったときは、その旨を当該加入者に通知しなければならない。
    \end{lawarticle}
  \end{lawstructure}
\end{fillblank}


\begin{answerbox}
  \begin{enumerate}[A: ]
    \item 資産管理運用機関等
    \item 個人別管理資産
  \end{enumerate}
\end{answerbox}



\begin{fillblank}[年金2 2018 問題1(3)D]
  \noindent \textbf{◯DB令}
  \begin{lawstructure}
    \begin{lawarticle}{第五十四条の九}{確定拠出年金又は独立行政法人勤労者退職金共済機構からの資産の移換の基準}
      法第八十二条の六第一項の政令で定める基準は、同項の移換又は引渡しを受けた額の算定の基礎となった期間の（　　A　　）を、厚生労働省令で定めるところにより、当該加入者に係る加入者期間に算入するものであることとする。
     \end{lawarticle}
  \end{lawstructure}
\end{fillblank}


\begin{answerbox}
  \begin{enumerate}[A: ]
    \item 全部又は一部
  \end{enumerate}
\end{answerbox}

\newpage

\subsection{中退共からDBへの資産の移換}

出題例なし

\newpage

